\honest{Not for the Sake of Happiness (Alone)}{Eliezer Yudkowsky}{2007}

Some years ago the futurist Greg Stock argued that pills would soon simulate the joy of scientific discovery, and be more fun than the real work: the real thing ``won't be able to compete.'' I told Stock: ``I agree that's possible, so I'll make sure never to take them.'' Stock seemed surprised, which surprised me.

Ethicists often argue ``as if all human desires are reducible, in principle, to the desire for ourselves and others to be happy.'' My example is Sam Harris, in a ``drive-by shooting.'' \nb{The passage of this kind found in Harris's \textsc{The End of Faith}, checked only in a quotation, says that ``questions of right and wrong are really questions about the happiness and suffering of sentient creatures.'' That is a claim about what matters morally, the kind philosophers call evaluative, not a claim about what we desire.} The question is not whether all happinesses share one scale, nor whether we can value only our own mental states, since we may care about others' happiness. It is whether we should care about the things that make us happy apart from the happiness they bring.

Laws against oral sex show moralists going astray, but one such case does not show that all values reduce to happiness. That we tend to do what makes us happy does not make happiness our only reason. First, if it were, we could care about others only as a means to our own warm glow. \nb{That tells against egoism. The view the post set out to examine counts others' happiness too, as the post itself allows.} Second, a result of an action need not be its only reason: if I take ibuprofen for a headache while writing, less pain is one consequence, not necessarily the main reason, and I can value something both for itself and as a means. Third, for all value to reduce to happiness, happiness must be the only thing that counts in every decision. I credit ``a Sober and Wilson paper, not sure which one.'' \nb{All three arguments are about what moves us; the book's account of how that bears on what we should care about comes later. The second is Joseph Butler's old objection to psychological egoism.}

Would I value art that no one ever sees, such as a screensaver in a closed room? ``I'd have to say no.'' Valuing a lifeless object for itself would be like valuing ice cream apart from anyone eating it. Everything I value involves people and their experiences somewhere. \nb{G. E. Moore argued the other way in 1903: a beautiful world that no one ever sees is better than ``one heap of filth.''} My intuition ``appears to require both the objective and subjective component to grant full value'': a real discovery, and a person's joy in it. I would be disturbed if people fell in love with mindless wallpaper in holodecks, even unknowingly, which matters if some agents could substitute zombies for people's loved ones without their knowledge. I would not enter a holodeck even if a pill made me forget it. \nb{This is Robert Nozick's ``experience machine'' of 1974, and Moore had described the same case. Hedonists answer that such intuitions are unreliable, and in one experiment people told they were already in the machine often chose to stay. The post does not consider this reply.} I also value freedom: an outcome counts for less if an unseen puppet master arranged it, even if no one knew, which matters if agents become powerful enough to tweak people's futures helpfully without their knowledge.

So my values are not strictly reducible to happiness. My decision system has ``a lot'' of terminal values, ``none of them strictly reducible to anything else'': art, science, love, lust, freedom, friendship. \nb{That rests on the linked ``Thou Art Godshatter''; this post argues only against reduction to happiness.} I value a life complicated enough to be challenging, with actual complications rather than the feeling of them, so being ``a pleasure center in a vat'' does not appeal to me. It would waste humanity's potential, which I want actually fulfilled.
